% The Logger's view of a request: a chain of self-describing events in the originating
% core's clock; consecutive events bound the stages.
\begin{tikzpicture}[font=\sffamily\scriptsize, >=Latex, node distance=3mm and 5mm,
  r/.style={draw=octrule, fill=octmem, rounded corners=2pt, minimum width=15mm, minimum height=7mm, font=\sffamily\footnotesize\bfseries},
  bus/.style={r, fill=octio, draw=octpurple},
  mem/.style={r, fill=white},
  w/.style={->, thick, octink}, e/.style={font=\ttfamily\scriptsize, text=octquiet}]
  \node[r] (cpu) {CPU};
  \node[r, right=of cpu] (l1) {L1};
  \node[bus, right=of l1] (rb) {REQ\_BUS};
  \node[r, right=of rb] (llc) {LLC};
  \node[bus, right=of llc] (sb) {RESP\_BUS};
  \node[r, right=of sb] (cpu2) {CPU};
  \node[bus, below=10mm of llc, xshift=-12mm] (mb1) {MEM\_BUS};
  \node[mem, right=of mb1] (dr) {DRAM};
  \node[bus, right=of dr] (mb2) {MEM\_BUS};
  \draw[w] (cpu) -- node[e, above] {ENTER} (l1);
  \draw[w] (l1) -- node[e, above] {ENTER} (rb);
  \draw[w] (rb) -- node[e, above] {EXIT} (llc);
  \draw[w] (llc) -- node[e, above] {ENTER} node[e, below] {SERVICE, EXIT} (sb);
  \draw[w] (sb) -- node[e, above] {EXIT} (cpu2);
  \node[e, above=0.5mm of cpu2] {EXIT $\rightarrow$ \code{finalizeEvents}};
  \draw[w, octquiet] (llc.south) |- node[e, below, pos=0.9] {miss} (mb1.north);
  \draw[w, octquiet] (mb1) -- node[e, above] {EXIT} (dr);
  \draw[w, octquiet] (dr) -- node[e, above] {ENTER, EXIT} (mb2);
  \draw[w, octquiet] (mb2.east) -- ++(4mm,0) |- node[e, above, pos=0.25] {EXIT} (llc.east);
  \node[text=octquiet, align=center, below=3mm of dr, text width=70mm] {every event records \code{(comp\_id, role, phase, cycle)}; the cycle is the originating core's clock whatever component stamps it, so stages are differences of named events and tile to Total};
\end{tikzpicture}
